% ============================================================
% Question Bank: ch07-07 Cylinder Surface Area and Volume
% Topic Title: Cylinder Surface Area and Volume
% CCSS: Supplementary (FL, VA)
% Questions: 27 (15 MC + 6 SA + 6 Graphical)
% ============================================================

% --- Multiple Choice Questions (15) ---

\begin{practiceQuestion}{7-7-q01}{mc}
\begin{questionText}
What is the volume of a cylinder with radius $5$ cm and height $6$ cm? Use $\pi \approx 3.14$.
\end{questionText}
\choiceA{$94.2$ cm$^3$}
\choiceB{$188.4$ cm$^3$}
\choiceC{$471$ cm$^3$}
\choiceD{$942$ cm$^3$}
\correctAnswer{C}
\explanation{The volume of a cylinder is $V = \pi r^2 h = 3.14 \times 25 \times 6 = 471$ cubic centimeters.}
\end{practiceQuestion}

\begin{practiceQuestion}{7-7-q02}{mc}
\begin{questionText}
Which formula gives the surface area of a cylinder?
\end{questionText}
\choiceA{$SA = 2\pi r^2 + \pi rh$}
\choiceB{$SA = 2\pi r^2 + 2\pi rh$}
\choiceC{$SA = \pi r^2 h$}
\choiceD{$SA = 2\pi r + 2\pi rh$}
\correctAnswer{B}
\explanation{The surface area of a cylinder is $SA = 2\pi r^2 + 2\pi rh$. The first term covers two circular bases and the second covers the lateral surface.}
\end{practiceQuestion}

\begin{practiceQuestion}{7-7-q03}{mc}
\begin{questionText}
A cylinder has a diameter of $10$ cm and a height of $7$ cm. What is its volume? Use $\pi \approx 3.14$.
\end{questionText}
\choiceA{$219.8$ cm$^3$}
\choiceB{$549.5$ cm$^3$}
\choiceC{$1{,}099$ cm$^3$}
\choiceD{$2{,}198$ cm$^3$}
\correctAnswer{B}
\explanation{First find the radius: $r = 10 \div 2 = 5$ cm. Then $V = \pi r^2 h = 3.14 \times 25 \times 7 = 549.5$ cubic centimeters.}
\end{practiceQuestion}

\begin{practiceQuestion}{7-7-q04}{mc}
\begin{questionText}
What is the surface area of a cylinder with radius $3$ cm and height $8$ cm? Use $\pi \approx 3.14$.
\end{questionText}
\choiceA{$150.72$ cm$^2$}
\choiceB{$207.24$ cm$^2$}
\choiceC{$263.76$ cm$^2$}
\choiceD{$414.48$ cm$^2$}
\correctAnswer{B}
\explanation{$SA = 2\pi r^2 + 2\pi rh = 2(3.14)(9) + 2(3.14)(3)(8) = 56.52 + 150.72 = 207.24$ square centimeters.}
\end{practiceQuestion}

\begin{practiceQuestion}{7-7-q05}{mc}
\begin{questionText}
The lateral surface area of a cylinder is given by $2\pi rh$. A tin can has radius $4$ cm and height $10$ cm. What is the lateral surface area? Use $\pi \approx 3.14$.
\end{questionText}
\choiceA{$125.6$ cm$^2$}
\choiceB{$251.2$ cm$^2$}
\choiceC{$502.4$ cm$^2$}
\choiceD{$1{,}004.8$ cm$^2$}
\correctAnswer{B}
\explanation{Lateral surface area: $2\pi rh = 2 \times 3.14 \times 4 \times 10 = 251.2$ square centimeters.}
\end{practiceQuestion}

\begin{practiceQuestion}{7-7-q06}{mc}
\begin{questionText}
A water bottle is cylindrical with radius $2.5$ cm and height $20$ cm. What is its volume? Use $\pi \approx 3.14$.
\end{questionText}
\choiceA{$157$ cm$^3$}
\choiceB{$392.5$ cm$^3$}
\choiceC{$785$ cm$^3$}
\choiceD{$196.25$ cm$^3$}
\correctAnswer{B}
\explanation{$V = \pi r^2 h = 3.14 \times 6.25 \times 20 = 392.5$ cubic centimeters.}
\end{practiceQuestion}

\begin{practiceQuestion}{7-7-q07}{mc}
\begin{questionText}
A cylinder has a volume of $628$ cm$^3$ and a radius of $5$ cm. Using $\pi \approx 3.14$, what is its height?
\end{questionText}
\choiceA{$4$ cm}
\choiceB{$8$ cm}
\choiceC{$16$ cm}
\choiceD{$25.12$ cm}
\correctAnswer{B}
\explanation{Rearrange $V = \pi r^2 h$ to get $h = V \div (\pi r^2) = 628 \div (3.14 \times 25) = 628 \div 78.5 = 8$ cm.}
\end{practiceQuestion}

\begin{practiceQuestion}{7-7-q08}{mc}
\begin{questionText}
A can of paint is cylindrical with diameter $12$ cm and height $14$ cm. What is the surface area? Use $\pi \approx 3.14$.
\end{questionText}
\choiceA{$527.52$ cm$^2$}
\choiceB{$665.28$ cm$^2$}
\choiceC{$753.6$ cm$^2$}
\choiceD{$1{,}507.2$ cm$^2$}
\correctAnswer{C}
\explanation{$r = 6$ cm. $SA = 2\pi r^2 + 2\pi rh = 2(3.14)(36) + 2(3.14)(6)(14) = 226.08 + 527.52 = 753.6$ square centimeters.}
\end{practiceQuestion}

\begin{practiceQuestion}{7-7-q09}{mc}
\begin{questionText}
When you unroll the lateral surface of a cylinder, what shape do you get?
\end{questionText}
\choiceA{Circle}
\choiceB{Triangle}
\choiceC{Rectangle}
\choiceD{Trapezoid}
\correctAnswer{C}
\explanation{The lateral surface of a cylinder unrolls into a rectangle. Its width equals the circumference ($2\pi r$) and its height equals the cylinder's height.}
\end{practiceQuestion}

\begin{practiceQuestion}{7-7-q10}{mc}
\begin{questionText}
A cylinder has radius $2$ cm and height $9$ cm. What is the volume? Use $\pi \approx 3.14$.
\end{questionText}
\choiceA{$56.52$ cm$^3$}
\choiceB{$113.04$ cm$^3$}
\choiceC{$226.08$ cm$^3$}
\choiceD{$452.16$ cm$^3$}
\correctAnswer{B}
\explanation{$V = \pi r^2 h = 3.14 \times 4 \times 9 = 113.04$ cubic centimeters.}
\end{practiceQuestion}

\begin{practiceQuestion}{7-7-q11}{mc}
\begin{questionText}
A cylindrical container has a volume of $942$ cm$^3$ and a height of $12$ cm. Using $\pi \approx 3.14$, what is the radius?
\end{questionText}
\choiceA{$5$ cm}
\choiceB{$10$ cm}
\choiceC{$15$ cm}
\choiceD{$25$ cm}
\correctAnswer{A}
\explanation{Rearrange $V = \pi r^2 h$: $r^2 = V \div (\pi h) = 942 \div (3.14 \times 12) = 942 \div 37.68 = 25$. So $r = 5$ cm.}
\end{practiceQuestion}

\begin{practiceQuestion}{7-7-q12}{mc}
\begin{questionText}
A cylindrical pillar has a diameter of $8$ m and a height of $3$ m. How much concrete is needed to fill it? Use $\pi \approx 3.14$.
\end{questionText}
\choiceA{$75.36$ m$^3$}
\choiceB{$150.72$ m$^3$}
\choiceC{$301.44$ m$^3$}
\choiceD{$602.88$ m$^3$}
\correctAnswer{B}
\explanation{$r = 4$ m. $V = \pi r^2 h = 3.14 \times 16 \times 3 = 150.72$ cubic meters.}
\end{practiceQuestion}

\begin{practiceQuestion}{7-7-q13}{mc}
\begin{questionText}
A net for a cylinder consists of which shapes?
\end{questionText}
\choiceA{Two circles and one rectangle}
\choiceB{One circle and one rectangle}
\choiceC{Two squares and one rectangle}
\choiceD{Two circles and one triangle}
\correctAnswer{A}
\explanation{A cylinder's net unfolds into two congruent circles (the bases) and one rectangle (the lateral surface whose width is the circumference).}
\end{practiceQuestion}

\begin{practiceQuestion}{7-7-q14}{mc}
\begin{questionText}
Doubling the radius of a cylinder while keeping the height the same causes the volume to:
\end{questionText}
\choiceA{double}
\choiceB{triple}
\choiceC{quadruple}
\choiceD{increase by $8$ times}
\correctAnswer{C}
\explanation{$V = \pi r^2 h$. If $r$ doubles to $2r$: $V_{new} = \pi (2r)^2 h = 4\pi r^2 h = 4V$. The volume quadruples.}
\end{practiceQuestion}

\begin{practiceQuestion}{7-7-q15}{mc}
\begin{questionText}
A cylindrical vase has a radius of $7$ cm and a height of $20$ cm. What is its volume? Use $\pi \approx \frac{22}{7}$.
\end{questionText}
\choiceA{$1{,}540$ cm$^3$}
\choiceB{$3{,}080$ cm$^3$}
\choiceC{$6{,}160$ cm$^3$}
\choiceD{$9{,}240$ cm$^3$}
\correctAnswer{B}
\explanation{$V = \pi r^2 h = \frac{22}{7} \times 49 \times 20 = 22 \times 7 \times 20 = 3{,}080$ cubic centimeters.}
\end{practiceQuestion}

% --- Short Answer Questions (6) ---

\begin{practiceQuestion}{7-7-q16}{sa}
\begin{questionText}
A cylindrical pipe has a radius of $6$ cm and a length of $25$ cm. What is the volume of the pipe? Use $\pi \approx 3.14$.
\end{questionText}
\correctAnswer{$2{,}826$ cm$^3$}
\explanation{$V = \pi r^2 h = 3.14 \times 36 \times 25 = 2{,}826$ cubic centimeters.}
\end{practiceQuestion}

\begin{practiceQuestion}{7-7-q17}{sa}
\begin{questionText}
A cylindrical candle has a diameter of $4$ cm and a height of $9$ cm. What is its surface area? Use $\pi \approx 3.14$.
\end{questionText}
\correctAnswer{$138.16$ cm$^2$}
\explanation{$r = 2$ cm. $SA = 2\pi r^2 + 2\pi rh = 2(3.14)(4) + 2(3.14)(2)(9) = 25.12 + 113.04 = 138.16$ square centimeters.}
\end{practiceQuestion}

\begin{practiceQuestion}{7-7-q18}{sa}
\begin{questionText}
A cylinder has a volume of $1{,}570$ cm$^3$ and a radius of $10$ cm. What is the height? Use $\pi \approx 3.14$.
\end{questionText}
\correctAnswer{$5$ cm}
\explanation{$h = V \div (\pi r^2) = 1{,}570 \div (3.14 \times 100) = 1{,}570 \div 314 = 5$ cm.}
\end{practiceQuestion}

\begin{practiceQuestion}{7-7-q19}{sa}
\begin{questionText}
A cylindrical water tank has a radius of $8$ m and a height of $5$ m. What is the lateral surface area? Use $\pi \approx 3.14$.
\end{questionText}
\correctAnswer{$251.2$ m$^2$}
\explanation{Lateral surface area: $2\pi rh = 2 \times 3.14 \times 8 \times 5 = 251.2$ square meters.}
\end{practiceQuestion}

\begin{practiceQuestion}{7-7-q20}{sa}
\begin{questionText}
A cylinder has the same radius and height, both equal to $4$ cm. What is its volume? Use $\pi \approx 3.14$.
\end{questionText}
\correctAnswer{$200.96$ cm$^3$}
\explanation{$V = \pi r^2 h = 3.14 \times 16 \times 4 = 200.96$ cubic centimeters.}
\end{practiceQuestion}

\begin{practiceQuestion}{7-7-q21}{sa}
\begin{questionText}
A cylindrical jar has a radius of $6$ cm and a height of $10$ cm. A label wraps around the lateral surface. What area does the label cover? Use $\pi \approx 3.14$.
\end{questionText}
\correctAnswer{$376.8$ cm$^2$}
\explanation{The label covers the lateral surface area: $2\pi rh = 2 \times 3.14 \times 6 \times 10 = 376.8$ square centimeters.}
\end{practiceQuestion}

% --- Graphical Questions (3 GMC + 3 GSA) ---

\begin{practiceQuestion}{7-7-q22}{gmc}
\begin{questionText}
The cylinder below has the dimensions shown. What is its volume? Use $\pi \approx 3.14$.

\begin{center}
\begin{tikzpicture}[scale=0.55]
  % Cylinder
  \draw[thick] (0,0) ellipse (2.5 and 0.75);
  \draw[thick] (-2.5,0) -- (-2.5,4.8);
  \draw[thick] (2.5,0) -- (2.5,4.8);
  \draw[thick] (0,4.8) ellipse (2.5 and 0.75);
  % Dimensions
  \draw[thick,funBlue] (0,4.8) -- (2.5,4.8);
  \node[funBlue, fill=white, inner sep=1pt] at (1.25,5.35) {$r = 4$};
  \draw[thick,funBlue,<->] (3.4,0) -- (3.4,4.8);
  \node[funBlue, fill=white, inner sep=1pt] at (4.35,2.4) {$h = 9$};
\end{tikzpicture}
\end{center}
\end{questionText}
\choiceA{$113.04$ cm$^3$}
\choiceB{$226.08$ cm$^3$}
\choiceC{$452.16$ cm$^3$}
\choiceD{$904.32$ cm$^3$}
\correctAnswer{C}
\explanation{$V = \pi r^2 h = 3.14 \times 16 \times 9 = 452.16$ cubic centimeters.}
\end{practiceQuestion}

\begin{practiceQuestion}{7-7-q23}{gmc}
\begin{questionText}
The cylinder below has a diameter of $14$ cm and a height of $10$ cm. What is the surface area? Use $\pi \approx 3.14$.

\begin{center}
\begin{tikzpicture}[scale=0.5]
  % Cylinder
  \draw[thick] (0,0) ellipse (2.7 and 0.78);
  \draw[thick] (-2.7,0) -- (-2.7,4.6);
  \draw[thick] (2.7,0) -- (2.7,4.6);
  \draw[thick] (0,4.6) ellipse (2.7 and 0.78);
  % Dimensions
  \draw[thick,funBlue,<->] (-2.7,5.25) -- (2.7,5.25);
  \node[funBlue, fill=white, inner sep=1pt] at (0,5.8) {$d = 14$};
  \draw[thick,funBlue,<->] (3.7,0) -- (3.7,4.6);
  \node[funBlue, fill=white, inner sep=1pt] at (4.8,2.3) {$h = 10$};
\end{tikzpicture}
\end{center}
\end{questionText}
\choiceA{$439.60$ cm$^2$}
\choiceB{$615.44$ cm$^2$}
\choiceC{$747.32$ cm$^2$}
\choiceD{$1{,}494.64$ cm$^2$}
\correctAnswer{C}
\explanation{$r = 7$ cm. $SA = 2\pi r^2 + 2\pi rh = 2(3.14)(49) + 2(3.14)(7)(10) = 307.72 + 439.6 = 747.32$ square centimeters.}
\end{practiceQuestion}

\begin{practiceQuestion}{7-7-q24}{gmc}
\begin{questionText}
The net of a cylinder is shown below. The two circles have radius $3$ cm and the rectangle has a height of $11$ cm. What is the total surface area? Use $\pi \approx 3.14$.

\begin{center}
\begin{tikzpicture}[scale=0.45]
  % Top circle
  \draw[thick,funBlue] (4.71,13.2) circle (1.5);
  \node[fill=white, inner sep=1pt] at (4.71,13.2) {\small $r = 3$};
  % Rectangle (lateral surface)
  \draw[thick,funBlue] (0,0) rectangle (9.42,11);
  \node[fill=white, inner sep=1pt] at (4.71,5.5) {\small $h = 11$};
  \node[fill=white, inner sep=1pt] at (4.71,-0.95) {\small width $= 2\pi r$};
  % Bottom circle
  \draw[thick,funBlue] (4.71,-3.2) circle (1.5);
  \node[fill=white, inner sep=1pt] at (4.71,-3.2) {\small $r = 3$};
\end{tikzpicture}
\end{center}
\end{questionText}
\choiceA{$207.24$ cm$^2$}
\choiceB{$263.76$ cm$^2$}
\choiceC{$320.28$ cm$^2$}
\choiceD{$527.52$ cm$^2$}
\correctAnswer{B}
\explanation{Two circles: $2 \times 3.14 \times 9 = 56.52$ cm$^2$. Rectangle (lateral): $2\pi r \times h = 2 \times 3.14 \times 3 \times 11 = 207.24$ cm$^2$. Total: $56.52 + 207.24 = 263.76$ square centimeters.}
\end{practiceQuestion}

\begin{practiceQuestion}{7-7-q25}{gsa}
\begin{questionText}
The cylinder below has the dimensions shown. Find its volume. Use $\pi \approx 3.14$.

\begin{center}
\begin{tikzpicture}[scale=0.56]
  % Cylinder
  \draw[thick] (0,0) ellipse (2.9 and 0.82);
  \draw[thick] (-2.9,0) -- (-2.9,3.4);
  \draw[thick] (2.9,0) -- (2.9,3.4);
  \draw[thick] (0,3.4) ellipse (2.9 and 0.82);
  % Dimensions
  \draw[thick,funBlue] (0,3.4) -- (2.9,3.4);
  \node[funBlue, fill=white, inner sep=1pt] at (1.45,4.0) {$r = 6$};
  \draw[thick,funBlue,<->] (3.9,0) -- (3.9,3.4);
  \node[funBlue, fill=white, inner sep=1pt] at (4.85,1.7) {$h = 5$};
\end{tikzpicture}
\end{center}
\end{questionText}
\correctAnswer{$565.2$ cm$^3$}
\explanation{$V = \pi r^2 h = 3.14 \times 36 \times 5 = 565.2$ cubic centimeters.}
\end{practiceQuestion}

\begin{practiceQuestion}{7-7-q26}{gsa}
\begin{questionText}
The cylinder below has a diameter of $8$ cm and a height of $12$ cm. Find its surface area. Use $\pi \approx 3.14$.

\begin{center}
\begin{tikzpicture}[scale=0.4]
  % Cylinder
  \draw[thick] (0,0) ellipse (2 and 0.6);
  \draw[thick] (-2,0) -- (-2,5);
  \draw[thick] (2,0) -- (2,5);
  \draw[thick] (0,5) ellipse (2 and 0.6);
  % Dimensions
  \draw[thick,funBlue,<->] (-2,5.8) -- (2,5.8) node[midway,above] {$d = 8$};
  \draw[thick,funBlue,<->] (2.8,0) -- (2.8,5) node[midway,right] {$h = 12$};
\end{tikzpicture}
\end{center}
\end{questionText}
\correctAnswer{$401.92$ cm$^2$}
\explanation{$r = 4$ cm. $SA = 2\pi r^2 + 2\pi rh = 2(3.14)(16) + 2(3.14)(4)(12) = 100.48 + 301.44 = 401.92$ square centimeters.}
\end{practiceQuestion}

\begin{practiceQuestion}{7-7-q27}{gsa}
\begin{questionText}
Two cylinders are shown below. Cylinder A has radius $2$ cm and height $15$ cm. Cylinder B has radius $5$ cm and height $4$ cm. Which cylinder has the greater volume, and by how much? Use $\pi \approx 3.14$.

\begin{center}
\begin{tikzpicture}[scale=0.46]
  % Cylinder A
  \draw[thick] (-4,0) ellipse (1.15 and 0.35);
  \draw[thick] (-5.15,0) -- (-5.15,6.5);
  \draw[thick] (-2.85,0) -- (-2.85,6.5);
  \draw[thick] (-4,6.5) ellipse (1.15 and 0.35);
  \node[below] at (-4,-0.7) {Cylinder A};
  \draw[thick,funBlue] (-4,6.5) -- (-2.85,6.5);
  \node[funBlue, fill=white, inner sep=1pt] at (-3.45,7.1) {\small $r=2$};
  \draw[thick,funBlue,<->] (-2.05,0) -- (-2.05,6.5);
  \node[funBlue, fill=white, inner sep=1pt] at (-1.0,3.25) {\small $h=15$};
  % Cylinder B
  \draw[thick] (6.2,0) ellipse (2.35 and 0.7);
  \draw[thick] (3.85,0) -- (3.85,2.45);
  \draw[thick] (8.55,0) -- (8.55,2.45);
  \draw[thick] (6.2,2.45) ellipse (2.35 and 0.7);
  \node[below] at (6.2,-1.0) {Cylinder B};
  \draw[thick,funBlue] (6.2,2.45) -- (8.55,2.45);
  \node[funBlue, fill=white, inner sep=1pt] at (7.35,3.1) {\small $r=5$};
  \draw[thick,funBlue,<->] (9.35,0) -- (9.35,2.45);
  \node[funBlue, fill=white, inner sep=1pt] at (10.35,1.225) {\small $h=4$};
\end{tikzpicture}
\end{center}
\end{questionText}
\correctAnswer{Cylinder B by $125.6$ cm$^3$}
\explanation{Cylinder A: $V = 3.14 \times 4 \times 15 = 188.4$ cm$^3$. Cylinder B: $V = 3.14 \times 25 \times 4 = 314$ cm$^3$. Cylinder B is greater by $314 - 188.4 = 125.6$ cubic centimeters.}
\end{practiceQuestion}
